% ============================================
% 第2章：为什么孩子会蹲在沙坑里搭三个小时的桥
% ============================================

\chapter{为什么孩子会蹲在沙坑里搭三个小时的桥}

\begin{classroom-scene}
\noindent \textbf{2026年5月22日，面向呼和浩特铁路局教师的线上分享会。}

\vspace{0.3cm}

\noindent 我在讲游戏化教学。开场前，我让大家扫描二维码，打开一个网页。那是一个桥梁建造游戏——SpanCraft的简化版。

\noindent 十分钟后，我注意到聊天区里，一个老师打了一行字：

\noindent\textbf{“我搭的桥又塌了。我再试一次。”}

\noindent 又过了五分钟，同一个人：

\noindent\textbf{“站住了！！！”}

\noindent 三个感叹号。一个成年人，一个铁路系统的工程师，在培训课上，因为一座虚拟的桥站住了，打了三个感叹号。

\noindent 这就是游戏。
\end{classroom-scene}

{\small\color{muted} [本章待撰写]}

\section{游戏的本质：主动选择困难}

\section{魔法圈：失败是安全的}

\section{胜任感：不是“我很厉害”，而是“我刚刚还很笨，现在不那么笨了”}

\section{八个维度：用游戏设计的眼光重新审视课堂}

\subsection{引人入胜（Engagement）}

\subsection{胜任感（Competence）}

\subsection{引导性（Onboarding）}

\subsection{美观性（Aesthetics）}

\subsection{故事性（Narrative）}

\subsection{挑战性（Challenge）}

\subsection{心流体验（Flow）}

\subsection{传播性（Shareability）}

\begin{shiyu-note}
{\small\color{muted} [师宇手记待补充]}
\end{shiyu-note}

\begin{decision-point}
{\small\color{muted} [教师决策时刻待补充]}
\end{decision-point}